% =============================================================
% §4 系统设计
% =============================================================
\section{系统设计}
\label{sec:system}

本系统由两个协同设计的组件构成：建模第~\ref{sec:intro} 节所列五项 AIGC 物理的推理仿真器（第~\ref{sec:simulator} 节），以及通过结构化状态与奖励整形将上述物理量暴露给策略的 AIGC 感知 PPO 调度器（第~\ref{sec:rl} 节）。

\subsection{AIGC 推理仿真器}
\label{sec:simulator}

仿真器以 Python 实现，采用 0.1~s 步长的离散事件驱动结构。每个 step 依次（a）回收已完成任务并释放资源，（b）按依赖与到达时间将任务从 \texttt{WAITING}~$\to$~\texttt{READY} 推进，（c）调用调度器，（d）启动被准入的任务并累积各服务器能耗。五项 AIGC 物理以 M1--M4 四个工程里程碑实现。

\textbf{M1：模型权重驻留。}每台服务器维护 \texttt{loaded\_models} 并附带引用计数；缺失模型触发未锁定模型的 LRU 驱逐与 \texttt{cold\_load\_sec} 延迟（LLaMA-7B/14~GB 为 5~s；LLaMA-70B-INT8/70~GB 为 25~s；按 NVMe 顺序读带宽标定）。

\textbf{M2：Prefill--Decode 两阶段任务与 KV cache。}每请求拆分为共享 \texttt{req\_id} 的 (Prefill, Decode) 对。KV 大小被编码进 \texttt{prefill.output\_size}（$\text{prompt\_tokens}\cdot\kappa_m/1024$~GB），使传输时间公式在 Decode 落到与 Prefill 不同的服务器时自动计入 \textbf{KV 迁移开销}；每个任务还携带 \texttt{kv\_cache\_GB} 字段，在执行期间占用显存。

\textbf{M3：连续批处理。}被准入的任务统计同模型同类型并行运行数 \texttt{batch\_size\_at\_admit}，其有效时间为 $T_{\text{solo}}(1+(\text{batch\_size}-1)\cdot\beta_k)$，其中 $\beta_{\text{decode}}{=}0.05$、$\beta_{\text{prefill}}{=}0.30$（与 vLLM 实测匹配）。\texttt{--ablation no\_batching} 开关同时关闭批处理\emph{并}强制每块 GPU 串行执行推理---模拟无迭代级批处理引擎的行为。

\textbf{M4：显存带宽下限 + 真实负载。}Decode 受显存带宽约束；我们以每模型 token 下限 $\phi_m$ 封顶有效时间（LLaMA-7B decode 20~ms/tok，按比例延展到 LLaMA-70B 的 80~ms/tok）。M4 还引入 Poisson 到达（\texttt{arrival\_time} 控制 \texttt{READY} 转换）与对齐 Azure LLM Inference Trace 2023 的对数正态 prompt/output 长度分布（prompt $\mu_{\log}{=}5.5,\sigma{=}1.0$；output $\mu_{\log}{=}4.5,\sigma{=}1.2$）。

\textbf{能耗模型。}每台服务器瞬时功率为 $P(t) = \text{idle\_W} + (\text{max\_W}-\text{idle\_W})\cdot u_s(t)$，分层系数对齐 NVIDIA TDP（云端 A100：50--400~W；边缘 A100：30--250~W；T4：20--70~W；Jetson AGX Orin：10--30~W）。系统每 token 能耗为积分功率除以已完成输出 token 数。

仿真器暴露单一调用的 \texttt{scheduler.schedule()} 接口，任何调度器---从 50 行的 RoundRobin 到我们的 PPO agent---都可即插即用，无需改动环境。图~\ref{fig:simulator-arch} 给出端到端数据流。

\begin{figure*}[t]
  \centering
  \includegraphics[width=0.70\textwidth]{trace_integration.pdf}
  \caption{AIGC 推理仿真器端到端架构。颜色带分别为：（紫）trace 源、（蓝）工作负载生成器、（红）含 M1--M4 物理模块的仿真环境、（黄）指标收集、（灰）输出文件。最小化的 \texttt{scheduler.schedule()} 接口将策略与环境解耦，使全部 11 个调度器与 12 组消融配置共用同一物理后端，保证对比公平。}
  \label{fig:simulator-arch}
\end{figure*}

\subsection{AIGC 感知强化学习调度器}
\label{sec:rl}

我们将单任务分派建模为 MDP，并用 PPO 求解；图~\ref{fig:rl-arch} 给出由此形成的决策闭环概览：从逐任务状态编码、经 actor-critic 策略，到 7 项奖励与闭合回路的 PPO 更新。

\begin{figure*}[t]
  \centering
  \includegraphics[width=0.66\textwidth]{rl_scheduler_arch.pdf}
  \caption{AIGC 感知 PPO 调度器的 MDP 决策闭环。对每个 \texttt{READY} 任务，状态编码器（\textcircled{1}）构建 96 维向量，其每服务器块携带五个 AIGC 特定特征（高亮）；actor-critic 策略（\textcircled{2}）输出带掩码的动作分布与价值估计；所选服务器被派发进仿真器（\textcircled{3}，五项 AIGC 物理）；7 项稠密奖励（\textcircled{4}）分为负载均衡、AIGC 感知与防坍缩三组；PPO 更新（\textcircled{5}）经策略梯度闭合回路。预训练正是让策略跳出 always-cloud 吸引子的关键。}
  \label{fig:rl-arch}
\end{figure*}

\textbf{状态。}对每个 \texttt{READY} 任务，编码器输出维度为 $(10+|\mathcal{M}|)+2+10N$ 的扁平向量：\emph{任务块}（计算需求、工作量、输入/输出大小、依赖计数、任务类型与模型 ID one-hot、归一化 KV 大小），2 维\emph{全局块}（队列长度、已完成比例），以及 $10$ 维的每服务器块，含五个通用特征（算力/显存利用率、带宽、队列长度、传输时间估计）与五个\textbf{AIGC 特定特征}：（i）已加载模型指示、（ii）归一化批占用率、（iii）剩余显存比例、（iv）归一化冷加载代价、（v）\emph{兄弟服务器指示}---对 Decode 任务，标记候选服务器是否就是 Prefill 运行的那台，使策略无需对请求 ID 进行推理即可直接读取 KV cache 本地性。当 edge${=}5$、$|\mathcal{M}|{=}4$ 时，状态维度为 96。

\textbf{动作。}在 $N$ 台服务器上离散选择，配以\textbf{动作掩码}（凡 \texttt{can\_allocate} 失败的服务器其概率清零）；掩码在每次分派时重算。

\textbf{奖励。}7 项稠密的逐决策奖励分量加权求和：
\begin{equation}
  \label{eq:reward}
  \begin{aligned}
    r =\;& 0.25\,r_{\text{time}}
        + 0.10\,r_{\text{balance}}
        + 0.10\,r_{\text{match}} \\
       &+ 0.15\,r_{\text{warm}}
        + 0.20\,r_{\text{batch}} \\
       &+ 0.10\,r_{\text{affinity}}
        + 0.10\,r_{\text{cloud}},
  \end{aligned}
\end{equation}
将三项负载均衡信号（$r_{\text{time}}$、$r_{\text{balance}}$、$r_{\text{match}}$ 分别惩罚长完成时间、服务器过载与算力不匹配）与四项 AIGC 特定信号相结合：$r_{\text{warm}}$（模型已加载时 $+1$，否则取 $-\text{cold\_load}/20$ 并 clip 至 $[-1,0]$）；$r_{\text{batch}}$ 鼓励加入同模型已有批（Decode 权重更高）；$r_{\text{affinity}}$ 在 Decode 与其 Prefill 同址时奖励 $+1$，在 KV 迁移成本非平凡时惩罚 $-0.5$；$r_{\text{cloud}}$（存在可行边缘替代时取 $-u_{\text{cloud}}$）抑制我们在早期实验中观察到的"always-cloud"吸引子（见 \S\ref{sec:implications}）。
典型权重将 0.45 分配给负载均衡组、0.45 分配给 AIGC 组、0.10 分配给防坍缩惩罚；组内各权重大致按各信号在诊断实验中的信息量比例设定。我们强调这一选择\emph{并非}精细调优：在从延迟优先到 AIGC 优先的四组备选权重向量下重新训练，两个目标在统计上均无变化（\S\ref{sec:weight-sens}）。

\textbf{网络结构与训练。}3 层多层感知机（MLP）编码器（$\text{state} \to 256 \to 128 \to 64$，LayerNorm + ReLU）接入 Actor 头（Linear $\to$ 带掩码 softmax）与 Critic 头（Linear $\to$ ReLU $\to$ Linear）。我们以 PPO~\cite{schulman2017ppo} 训练（clip $\epsilon{=}0.2$）、GAE~\cite{schulman2015gae}（$\gamma{=}0.95$、$\lambda{=}0.90$）与熵正则（系数 $0.05$，预训练后降至 $0.01$）；每次更新使用 32 条缓存 transition、4 个 epoch。关键地，我们在在线训练前进行 \textbf{20 个 episode 的预训练}（约 10--15~s 墙钟），起手 30 步使用均匀随机动作 warm-up 初始化；若无此预训练，策略会坍缩到"always-cloud"，\slo{} 下降 $38\%$（\S\ref{sec:ablation}）。每个 AIGC 组件由构造器开关控制，从而支持 \S\ref{sec:ablation} 的 12 组件消融。
